% Unit 063 English translation of chapter4.tex lines 2500--2745.
% Source: Methods of Algebra, Volume 2, commit
% 9a5803ff77dd3257484cb177f851a73770a59dd3 (CC BY 4.0).
% stable-unit-id: o014.aljabr2.chapter4.k-flat-derived-tensor-b
% Coverage: end of Section 4.12 and all Chapter 4 exercises (17 exercises and
% 10 hints); closes chapter4.tex exactly at the end of the Exercises environment.

% segment-id: o014.li2.chapter4.k-flat-derived-tensor-b.co001
\hypertarget{o014.aljabr2.chapter4.k-flat-derived-tensor-b}{ }
\begin{corollary}\label{prop:otimesL-Tor}
	\index[sym1]{Tor}
	The functor $\otimesL[R]$ of Definition~\ref{def:otimesL}, under the
	assumption that $A$ or $B$ is flat as a $\Bbbk$-module, satisfies the
	canonical isomorphism
	\[ \Hm^{-n}\left( X \otimesL[R] Y \right) \simeq \Tor^R_n(X, Y), \quad n \in \Z, \]
	where $X$ is an $(A,R)$-bimodule, $Y$ is an $(R,B)$-bimodule, and
	$\Tor^R_n$ is as in Definition--Proposition~\ref{def:Tor-classical}.
\end{corollary}
% segment-id: o014.li2.chapter4.k-flat-derived-tensor-b.pf001
\begin{proof}
	The left-hand side is obtained by composing the diagram of
	Proposition~\ref{prop:otimesL-enrichment} in the direction
% segment-id: o014.li2.chapter4.k-flat-derived-tensor-b.g001
	\begin{tikzpicture}[xscale=0.3, yscale=0.3, baseline=(O)]
		\draw[->] (0,0) -- (1,0) -- (1, -1);
		\coordinate (O) at (0, -0.7);
	\end{tikzpicture}
	whereas the right-hand side is obtained by composing it in the direction
% segment-id: o014.li2.chapter4.k-flat-derived-tensor-b.g002
	\begin{tikzpicture}[xscale=0.3, yscale=0.3, baseline=(O)]
		\draw[->] (0,0) -- (0, -1) -- (1, -1);
		\coordinate (O) at (0, -0.7);
	\end{tikzpicture}.
\end{proof}

% segment-id: o014.li2.chapter4.k-flat-derived-tensor-b.pr001
\begin{proposition}[Associativity constraint]\label{prop:otimesL-associativity-constraint}
	\index{derived tensor product!associativity constraint}
	Let $A$, $B$, $R$, and $S$ be $\Bbbk$-algebras. Consider the following two
	functors:
	\begin{gather*}
		\cate{D}((A,R)\dcate{Mod}) \times \cate{D}((R,S)\dcate{Mod}) \times \cate{D}((S,B)\dcate{Mod}) \to \cate{D}((A,B)\dcate{Mod}), \\
		(X, Y, Z) \mapsto (X \otimesL[R] Y) \otimesL[S] Z, \\
		(X, Y, Z) \mapsto X \otimesL[R] (Y \otimesL[S] Z).
	\end{gather*}
	If $A$ and $B$ are both flat as $\Bbbk$-modules, there is a canonical
	isomorphism
	\[ a(X,Y,Z): (X \otimesL[R] Y) \otimesL[S] Z \rightiso X \otimesL[R] (Y \otimesL[S] Z). \]
\end{proposition}
% segment-id: o014.li2.chapter4.k-flat-derived-tensor-b.pf002
\begin{proof}
	Fix $Y$. By Proposition~\ref{prop:K-flat}, we may assume that $X_R$ and
	${}_R Z$ are K-flat. Then every $\otimesL[R]$ can be computed at the level
	of complexes, and the assertion reduces to the associativity constraint for
	tensor products \cite[Proposition 6.5.12]{Li1}.
\end{proof}

% segment-id: o014.li2.chapter4.k-flat-derived-tensor-b.p001
Under these flatness assumptions, taking $\otimesL$ of three or more objects
in different orders yields results that are mutually isomorphic through the
associativity constraint. These isomorphisms also satisfy coherence, which can
be formulated by imitating the pentagon axiom for monoidal categories; see
\cite[Definition 3.1.1]{Li1}. Again, one takes K-flat resolutions and checks
the assertion at the level of complexes. To keep the discussion brief, this
book presents only the special case of a commutative ring.

% segment-id: o014.li2.chapter4.k-flat-derived-tensor-b.e001
\begin{example}[$\otimesL$ over a commutative ring]\label{eg:commutative-otimesL}
	\index{derived tensor product!commutativity constraint}
	Let $R$ be a commutative ring. In the definition
	$\otimesL:=\otimesL[R]$, take every ring to be $R$, in particular
	$\Bbbk=R$. Then $(R,R)\dcate{Mod}=R\dcate{Mod}$, and all the preceding
	flatness assumptions are automatically satisfied. Consequently,
	$\Hm^{-n}(X\otimesL Y)\simeq\Tor^R_n(X,Y)$ for all $R$-modules $X$ and
	$Y$.

	This makes $\cate{D}(R\dcate{Mod})$ a monoidal category with respect to
	$\otimesL$ in the sense of \cite[\S 3.1]{Li1}. Its unit object is $R$
	itself, and its associativity constraint comes from
	$a=(a(X,Y,Z))_{X,Y,Z}$ in
	Proposition~\ref{prop:otimesL-associativity-constraint}. Since $R$ is a flat
	$R$-module, it is K-flat (Proposition~\ref{prop:flat-K-flat}); hence both
	$R\otimesL(\cdot)$ and $(\cdot)\otimesL R$ can be computed at the level of
	complexes, making all the monoidal-category axioms easy to prove.
	Furthermore, $\cate{D}(R\dcate{Mod})$ is a symmetric monoidal category in
	the sense of \cite[Definition 3.3.1]{Li1}. Its braiding (or commutativity
	constraint) $c(X,Y):X\otimesL Y\rightiso Y\otimesL X$ comes from the
	commutativity constraint at the level of complexes, and its axioms are also
	checked there. In view of Proposition~\ref{prop:double-cplx-swap}, swapping
	the superscripts in the total complex causes no difficulty here.
\end{example}

% segment-id: o014.li2.chapter4.k-flat-derived-tensor-b.e002
\begin{example}[$\Tor$-algebra]
	Continue with the setting of Example~\ref{eg:commutative-otimesL}, with $R$
	commutative. Identify the complexes of $R$-modules $X,Y,Z,W$ with their
	images in the derived category. For every $(p,q)\in\Z^2$, there is a
	canonical morphism
% segment-id: o014.li2.chapter4.k-flat-derived-tensor-b.q001
	\begin{multline*}
		\Hm^{-p}\left( X \otimesL Y \right) \otimes \Hm^{-q}\left( Z \otimesL W \right) \to \Hm^{-p-q}\left( (X \otimesL Y) \otimesL (Z \otimesL W) \right) \\
		\simeq \Hm^{-p-q}\left( (X \otimesL Z) \otimesL (Y \otimesL W) \right) \to \Hm^{-p-q} \left( (X \otimes Z) \otimesL (Y \otimes W) \right),
	\end{multline*}
	where the first part applies Proposition~\ref{prop:bifunctor-cup} (take
	$F=\otimes$), the second uses the associativity and commutativity constraints
	on $(\cate{D}(R\dcate{Mod}),\otimesL)$, and the third comes from the
	canonical morphism accompanying $\otimesL$ as a left derived bifunctor:
% segment-id: o014.li2.chapter4.k-flat-derived-tensor-b.q002
	\[ QX_1 \otimesL QX_2 \to Q(X_1 \otimes X_2), \quad X_1, X_2 \in \Obj(\cate{K}(R\dcate{Mod})), \]
	For precision, the localization functor
	$Q:\cate{K}(R\dcate{Mod})\to\cate{D}(R\dcate{Mod})$ has been written
	explicitly again here.
% segment-id: o014.li2.chapter4.k-flat-derived-tensor-b.l001
	\begin{itemize}
% segment-id: o014.li2.chapter4.k-flat-derived-tensor-b.i001
		\item The composite of the three canonical morphisms above gives an
		``external product'' for the $\Tor$ functor:
		\[ \Tor^R_p(X, Y) \otimes \Tor^R_q(Z, W) \to \Tor^R_{p+q}(X \otimes Z, Y \otimes W). \]
% segment-id: o014.li2.chapter4.k-flat-derived-tensor-b.i002
		\item Take $Z=X$ and $W=Y$ to be $R$-algebras (concentrated in degree
		$0$). Their multiplications give homomorphisms of $R$-modules
		$X\otimes X\to X$ and $Y\otimes Y\to Y$. Composing the external product
		with these homomorphisms gives
		\[ \Tor^R_p(X, Y) \otimes \Tor^R_q(X, Y) \to \Tor^R_{p+q}(X,Y), \]
		called the ``internal product'' for the $\Tor$ functor. This makes
		$\bigoplus_{n \in \Z} \Tor^R_n(X, Y)$ a $\Z$-graded $R$-algebra in the
		sense of \cite[Definition 7.4.1]{Li1}, called the \emph{$\Tor$-algebra}.
		It has nonzero components only in nonnegative degrees, and its degree-zero
		component is the subalgebra $\Tor^R_0(X,Y)=X\otimes Y$, whose definition
		may be found in \cite[Definition 7.3.1]{Li1}. The associativity required
		for the product can be verified by reducing it to the associativity
		constraint for $\otimesL$ and to the associativity of multiplication on
		$X$ and $Y$ separately.
		\index{Tor-algebra@$\Tor$-algebra}
% segment-id: o014.li2.chapter4.k-flat-derived-tensor-b.i003
		\item Suppose further that $X$ and $Y$ are commutative $R$-algebras. In
		this case, the $\Tor$-algebra is also an ``anticommutative''
		$\Z$-graded algebra in the sense of
		\cite[Definition 7.4.3]{Li1}\footnote{The literal term is somewhat
		misleading, because this property is a natural manifestation of
		commutativity. Graded-commutative is more appropriate, as in
		Example~\ref{eg:graded-module}.}: for all $p,q\in\Z$,
		\[ a \in \Tor^R_p(X, Y), \; b \in \Tor^R_q(X, Y) \implies ab = (-1)^{pq} ba. \]
		Since both $X\otimes X\to X$ and $Y\otimes Y\to Y$ are commutative,
		where does the sign come from? Its source is the definition of the
		external product. Let $C_1$ and $C_2$ be two copies of $X\otimesL Y$.
		Then $ab$ and $ba$ differ through the isomorphism
		$C_1\otimesL C_2\rightiso C_2\otimesL C_1$, namely the commutativity
		constraint for $\otimesL$. If $C_1$ and $C_2$ are represented by K-flat
		complexes, this isomorphism of total complexes is exactly the isomorphism
		$r_{C_1^\bullet \otimes C_2^\bullet}$ of
		Proposition~\ref{prop:double-cplx-swap}; on the component of bidegree
		$(p,q)$, it produces the sign $(-1)^{pq}$.
	\end{itemize}
\end{example}

% segment-id: o014.li2.chapter4.k-flat-derived-tensor-b.p002
Return to general $\Bbbk$-algebras. To discuss the adjunction between
$\otimesL$ and $\RHom$, one more lemma is needed. Below,
$\RHom_{(A,B)}$ denotes $\RHom$ in $(A,B)\dcate{Mod}$; analogous notation is
used for other categories of bimodules (or of left or right modules).

% segment-id: o014.li2.chapter4.k-flat-derived-tensor-b.le001
\begin{lemma}\label{prop:RHom-enrichment}
	Fix $\Bbbk$-algebras $A$, $B$, and $R$. Suppose that $B$ is flat as a
	$\Bbbk$-module. Then the forgetful functor from $(A,B)\dcate{Mod}$ to
	$A\dcate{Mod}$ preserves K-injective complexes, and $\RHom_A$ has a
	canonical lift to a functor
	\[ \cate{D}\left((A,R)\dcate{Mod}\right)^{\opp} \times \cate{D}\left((A,B)\dcate{Mod}\right) \to \cate{D}\left((R,B)\dcate{Mod}\right), \]
	in other words, there is a canonical $2$-cell diagram
% segment-id: o014.li2.chapter4.k-flat-derived-tensor-b.g003
	\[\begin{tikzcd}
		\cate{D}\left((A,R)\dcate{Mod}\right)^{\opp} \times \cate{D}\left((A,B)\dcate{Mod}\right) \arrow[r] \arrow[d, "\text{forgetful}"'] & \cate{D}\left((R,B)\dcate{Mod}\right) \arrow[d, "\text{forgetful}"] \arrow[Rightarrow, ld, "\sim", sloped] \\
		\cate{D}(A\dcate{Mod})^{\opp} \times \cate{D}(A\dcate{Mod}) \arrow[r, "{\RHom_A}"'] & \cate{D}(\cate{Ab}) .
	\end{tikzcd}\]

	Likewise, if $A$ is flat as a $\Bbbk$-module, then $\RHom_B$ has an
	analogous lift.
\end{lemma}
% segment-id: o014.li2.chapter4.k-flat-derived-tensor-b.pf003
\begin{proof}
	Because $B$ is a flat $\Bbbk$-module, the functor $Y\mapsto{}_A Y$ has the
	exact left adjoint $(\cdot)\dotimes{\Bbbk}B$, and therefore preserves
	K-injective complexes.

	Next take $X\in\Obj(\cate{K}((A,R)\dcate{Mod}))$ and
	$Y\in\Obj(\cate{K}(A,B)\dcate{Mod})$, with $Y$ K-injective. By the
	preceding step, $\Hom^\bullet\left({}_A X,{}_A Y\right)$ computes
	$\RHom_A({}_A X,{}_A Y)$, while at the same time it is a complex of
	$(R,B)$-bimodules. This completes the lift.
\end{proof}

% segment-id: o014.li2.chapter4.k-flat-derived-tensor-b.th001
\begin{theorem}[Adjunction between $\RHom$ and $\otimesL$]\label{prop:RHom-otimesL-adjunction}
	\index{$\RHom$ functor}
	\index{derived tensor product}
	Let $A$, $B$, and $R$ be $\Bbbk$-algebras. Below, $X$, $Y$, and $Z$ denote
	objects of $\cate{D}((A,R)\dcate{Mod})$,
	$\cate{D}((R,B)\dcate{Mod})$, and
	$\cate{D}((A,B)\dcate{Mod})$, respectively. Suppose that $B$ is flat as a
	$\Bbbk$-module. Then there is a canonical isomorphism in
	$\cate{D}(\Bbbk\dcate{Mod})$
% segment-id: o014.li2.chapter4.k-flat-derived-tensor-b.q003
	\begin{equation*}
		\RHom_{(A,B)} \left( X \otimesL[R] Y, Z \right) \rightiso \RHom_{(R,B)}\left( Y, \RHom_A \left( {}_A X, {}_A Z \right) \right).
	\end{equation*}

	If instead $A$ is assumed flat as a $\Bbbk$-module, there is a canonical
	isomorphism in $\cate{D}(\Bbbk\dcate{Mod})$
% segment-id: o014.li2.chapter4.k-flat-derived-tensor-b.q004
	\begin{equation*}
		\RHom_{(A,B)}\left( X \otimesL[R] Y, Z \right) \rightiso \RHom_{(A,R)}\left( X, \RHom_B\left( Y_B , Z_B \right) \right).
	\end{equation*}
\end{theorem}
% segment-id: o014.li2.chapter4.k-flat-derived-tensor-b.pf004
\begin{proof}
	Consider the first equality. Suppose that $Y$ is K-projective and $Z$ is
	K-injective. Lemma~\ref{prop:K-projective-flat} ensures that ${}_R Y$ is
	K-flat, while Lemma~\ref{prop:RHom-enrichment} ensures that ${}_A Z$
	remains K-injective. Every displayed $\RHom$ and $\otimesL[R]$ can therefore
	be computed at the level of complexes. The desired isomorphism then follows
	from Proposition~\ref{prop:tensor-Hom-cplx}.
\end{proof}

% segment-id: o014.li2.chapter4.k-flat-derived-tensor-b.e003
\begin{example}[Adjunction for change of rings]\label{eg:change-of-ring-adjunction}
	In the setting of Example~\ref{eg:change-of-rings} (take a homomorphism
	$R\to S$ of $\Bbbk$-algebras, $A=S$, $B=\Bbbk$, and $X=S$), the first
	equality in Theorem~\ref{prop:RHom-otimesL-adjunction} becomes the canonical
	isomorphism
% segment-id: o014.li2.chapter4.k-flat-derived-tensor-b.q005
	\begin{gather*}
		\RHom_S\left( \mathrm{L} {}_{R \to S} P(Y), Z \right) \simeq \RHom_R\left( Y, {}_R Z \right), \\
		Y \in \Obj(\cate{D}(R\dcate{Mod})), \quad Z \in \Obj(\cate{D}(S\dcate{Mod})).
	\end{gather*}
	The only point still requiring explanation is
	$\RHom_S({}_S S,{}_S Z)\simeq{}_R Z$, which follows immediately from the
	isomorphism at the level of complexes
% segment-id: o014.li2.chapter4.k-flat-derived-tensor-b.q006
	\[ \Hom^\bullet\left({}_S S, \cdot\right) \simeq {}_R (\cdot) : \cate{C}(S\dcate{Mod}) \to \cate{C}(R\dcate{Mod}). \]

	If, furthermore, both $R$ and $S$ are commutative rings and $\Bbbk:=R$,
	then $\RHom_S$ (respectively, $\RHom_R$) takes values in
	$\cate{D}(S\dcate{Mod})$ (respectively, $\cate{D}(R\dcate{Mod})$). The
	adjunction isomorphism can then be written as the canonical isomorphism in
	$\cate{D}(R\dcate{Mod})$:
% segment-id: o014.li2.chapter4.k-flat-derived-tensor-b.q007
	\[ {}_R \RHom_S\left( \mathrm{L} {}_{R \to S} P(Y), Z \right) \simeq \RHom_R\left( Y, {}_R Z \right) . \]
	Everything can be verified at the level of complexes by taking K-injective
	and K-projective resolutions. The case of right modules is similar.
\end{example}


% segment-id: o014.li2.chapter4.k-flat-derived-tensor-b.x001
\begin{Exercises}
	% Reference: [KS06, Exercise 10.10]
% segment-id: o014.li2.chapter4.k-flat-derived-tensor-b.ex001
	\item Let $(\mathcal{D}, T, \mathcal{H})$ be a pretriangulated category
	(respectively, a triangulated category). Define a new family of triangles
	$\mathcal{H}^-$ by
% segment-id: o014.li2.chapter4.k-flat-derived-tensor-b.q008
	\[ [X \xrightarrow{f} Y \xrightarrow{g} Z \xrightarrow{h} TX] \in \mathcal{H} \iff [X \xrightarrow{f} Y \xrightarrow{g} Z \xrightarrow{-h} TX] \in \mathcal{H}^- . \]
% segment-id: o014.li2.chapter4.k-flat-derived-tensor-b.l002
	\begin{enumerate}[(i)]
% segment-id: o014.li2.chapter4.k-flat-derived-tensor-b.i004
		\item Prove that $(\mathcal{D},T,\mathcal{H}^-)$ is also a
		pretriangulated category (respectively, a triangulated category).
% segment-id: o014.li2.chapter4.k-flat-derived-tensor-b.i005
		\item Prove that $(\mathcal{D},T,\mathcal{H})$ and
		$(\mathcal{D},T,\mathcal{H}^-)$ are equivalent as pretriangulated
		categories.
	\end{enumerate}

% segment-id: o014.li2.chapter4.k-flat-derived-tensor-b.ex002
	\item Let $(\mathcal{D},T,\mathcal{H})$ be a pretriangulated category.
	Prove that every monomorphism (respectively, epimorphism) $u:X\to Y$ in
	$\mathcal{D}$ has a left inverse (respectively, a right inverse). Use this
	to show that a pretriangulated category that is also abelian must be split
	(Definition~\ref{def:semisimple}).

% segment-id: o014.li2.chapter4.k-flat-derived-tensor-b.hn001
	\begin{hint}
		Place the monomorphism $u$ in a distinguished triangle
		$X\xrightarrow{u}Y\xrightarrow{v}Z\xrightarrow{w}TX$. From
		$u\circ T^{-1}w=0$, deduce that $w=0$, and then apply
		Proposition~\ref{prop:dt-sum-1}.
	\end{hint}

% segment-id: o014.li2.chapter4.k-flat-derived-tensor-b.ex003
	\item Apply the preceding exercise to give an example of an abelian category
	$\mathcal{A}$ for which neither $\cate{K}(\mathcal{A})$ nor
	$\cate{D}(\mathcal{A})$ is an abelian category.

% segment-id: o014.li2.chapter4.k-flat-derived-tensor-b.hn002
	\begin{hint}
		Consider the morphism $f:\Z/p^2\Z\twoheadrightarrow\Z/p\Z$ in
		$\mathcal{A}=\cate{Ab}$, where $p$ is prime. If
		$\cate{K}(\mathcal{A})$ were abelian, consider the epi--mono
		factorizations of $\cate{K}f$ and $\Hm^0(\cate{K}f)$, then express both
		as projections onto direct summands and inclusion morphisms to obtain a
		contradiction.
	\end{hint}

% segment-id: o014.li2.chapter4.k-flat-derived-tensor-b.ex004
	\item Let $\mathcal{A}$ be an abelian category. Prove that
	$\cate{D}(\mathcal{A})$ is an abelian category if and only if
	$\mathcal{A}$ is split.

% segment-id: o014.li2.chapter4.k-flat-derived-tensor-b.hn003
	\begin{hint}
		For the ``only if'' direction, embed $\mathcal{A}$ in
		$\cate{D}(\mathcal{A})$ and apply the preceding exercises. For the ``if''
		direction, show that when $\mathcal{A}$ is split, taking cohomology gives
		an equivalence from $\cate{D}(\mathcal{A})$ to $\mathcal{A}^{\Z}$.
	\end{hint}

% segment-id: o014.li2.chapter4.k-flat-derived-tensor-b.ex005
	\item It is often said that the octahedral axiom (TR5) for a triangulated
	category replaces the isomorphism theorem
	$\dfrac{X/Z}{Y/Z}\simeq X/Y$ in an abelian category, where
	$X\supset Y\supset Z$ (Theorem~\ref{prop:Abel-cat-isom-thm} (ii)). Explain
	explicitly what this statement means.

% segment-id: o014.li2.chapter4.k-flat-derived-tensor-b.ex006
	\item ($\mathrm{K}_0$ of a pretriangulated category) For every
	pretriangulated category $(\mathcal{D},T,\mathcal{H})$, define
	$\mathrm{K}_0(\mathcal{D})$ to be the quotient of the free $\Z$-module
	generated by $\Obj(\mathcal{D})$ by the relations
% segment-id: o014.li2.chapter4.k-flat-derived-tensor-b.q009
	\[ \text{there is a distinguished triangle}\; X \to Y \to Z \xrightarrow{+1} \implies [X] - [Y] + [Z] = 0, \]
	where $[X]\in\mathrm{K}_0(\mathcal{D})$ denotes the equivalence class
	containing $X\in\Obj(\mathcal{D})$.
% segment-id: o014.li2.chapter4.k-flat-derived-tensor-b.l003
	\begin{enumerate}[(i)]
% segment-id: o014.li2.chapter4.k-flat-derived-tensor-b.i006
		\item Prove that $[0]=0$, $[TX]=-[X]$, and
		$[X\oplus Y]=[X]+[Y]$.
% segment-id: o014.li2.chapter4.k-flat-derived-tensor-b.i007
		\item Show that every triangulated functor
		$F:\mathcal{D}\to\mathcal{D}'$ induces a canonical homomorphism
		$\mathrm{K}_0(\mathcal{D})\to\mathrm{K}_0(\mathcal{D}')$.
% segment-id: o014.li2.chapter4.k-flat-derived-tensor-b.i008
		\item Consider a cohomological functor $H:\mathcal{D}\to\mathcal{A}$,
		where $\mathcal{A}$ is an abelian category. Suppose that, for every $X$,
		$|n|\gg_X0\implies H(T^nX)=0$. Define a natural homomorphism
		$\mathrm{K}_0(\mathcal{D})\to\mathrm{K}_0(\mathcal{A})$.

% segment-id: o014.li2.chapter4.k-flat-derived-tensor-b.hn004
		\begin{hint}
			Send $[X]$ to $\sum_n(-1)^n[H(T^nX)]$. Apply
			Lemma~\ref{prop:K0-prep} (iv).
		\end{hint}
% segment-id: o014.li2.chapter4.k-flat-derived-tensor-b.i009
		\item Consider the special case
		$\mathcal{D}:=\cate{D}^{\bdd}(\mathcal{A})$. Give mutually inverse
		homomorphisms
% segment-id: o014.li2.chapter4.k-flat-derived-tensor-b.g004
		\[\begin{tikzcd}[row sep=tiny]
			\mathrm{K}_0(\mathcal{A}) \arrow[shift left, r] & \mathrm{K}_0(\cate{D}^{\bdd}(\mathcal{A})) \arrow[shift left, l] \\
			{[A]} \arrow[mapsto, r] & {[A]} \\
			\sum_n (-1)^n {[\Hm^n(X)]} & {[X]} \arrow[mapsto, l]
		\end{tikzcd}\]
		where $A\in\Obj(\mathcal{A})$, also regarded as a complex concentrated in
		degree $0$.
% segment-id: o014.li2.chapter4.k-flat-derived-tensor-b.hn005
		\begin{hint}
			One composite is plainly the identity; the other requires truncating the
			complex.
		\end{hint}
	\end{enumerate}
	\index{K0 group@$\mathrm{K}_0$ group} \index[sym1]{K0}

% segment-id: o014.li2.chapter4.k-flat-derived-tensor-b.ex007
	\item Let $X\in\Obj(\cate{D}(\mathcal{A}))$. Prove that
% segment-id: o014.li2.chapter4.k-flat-derived-tensor-b.q010
	\begin{gather*}
		X \in \cate{D}^{\leq 0}(\mathcal{A}) \iff \forall Z \in \Obj(\cate{D}^{\geq 1}(\mathcal{A})), \; \Hom_{\cate{D}(\mathcal{A})}(X, Z) = 0, \\
		X \in \cate{D}^{\geq 0}(\mathcal{A}) \iff \forall Z \in \Obj(\cate{D}^{\leq -1}(\mathcal{A})), \; \Hom_{\cate{D}(\mathcal{A})}(Z, X) = 0.
	\end{gather*}
% segment-id: o014.li2.chapter4.k-flat-derived-tensor-b.hn006
	\begin{hint}
		Apply Proposition~\ref{prop:derived-cat-orthogonality} (i), together with
		the adjoint pairs and distinguished triangle in
		Proposition~\ref{prop:derived-truncation-adjunction}.
	\end{hint}

% segment-id: o014.li2.chapter4.k-flat-derived-tensor-b.ex008
	\item In the setting of
	Theorem~\ref{prop:localization-triangulated-composite}, suppose that $F'$
	satisfies $F'(\Obj(\mathcal{N}'))\subset\Obj(\mathcal{N}'')$ (regard this
	as a kind of ``exactness''). If $\mathcal{D}$ has an $F$-injective or an
	$F$-projective subcategory, prove that the canonical morphism
	\[ \mathrm{R}^{\mathcal{N}''}_{\mathcal{N}} (F' F) \to \left( \mathrm{R}^{\mathcal{N}''}_{\mathcal{N}'} F' \right) \left( \mathrm{R}^{\mathcal{N}'}_{\mathcal{N}} F \right) \]
	or
	\[ \left( \mathrm{L}^{\mathcal{N}''}_{\mathcal{N'}} F' \right) \left( \mathrm{L}^{\mathcal{N}'}_{\mathcal{N}} F \right) \to \mathrm{L}^{\mathcal{N}''}_{\mathcal{N}} (F' F) \]
	is an isomorphism.

% segment-id: o014.li2.chapter4.k-flat-derived-tensor-b.ex009
	\item In the setting of
	Theorem~\ref{prop:localization-triangulated-composite}, show that the
	following diagrams of functors and morphisms between them commute (the
	notation $\mathcal{N}$ and so forth is suppressed):
% segment-id: o014.li2.chapter4.k-flat-derived-tensor-b.g005
	\[\begin{tikzcd}
		\mathrm{R}(F'' F' F) \arrow[r] \arrow[d] & \mathrm{R}(F'' F') \mathrm{R}F \arrow[d] \\
		\mathrm{R}F'' \mathrm{R}(F'F) \arrow[r] & \mathrm{R}F'' \mathrm{R}F' \mathrm{R}F
	\end{tikzcd} \quad \begin{tikzcd}
		\mathrm{L}(F'' F' F) & \mathrm{L}(F'' F') \mathrm{L}F \arrow[l] \\
		\mathrm{L}F'' \mathrm{L}(F'F) \arrow[u] & \mathrm{L}F'' \mathrm{L}F' \mathrm{L}F \arrow[u] \arrow[l] .
	\end{tikzcd}\]

% segment-id: o014.li2.chapter4.k-flat-derived-tensor-b.ex010
	\item For $F=\Hom$, express explicitly the canonical morphism of
	Proposition~\ref{prop:bifunctor-cup} in the form
	$\Ext^{q-p}(X,Y)\to\Hom\left(\Hm^p(X),\Hm^q(Y)\right)$.

% segment-id: o014.li2.chapter4.k-flat-derived-tensor-b.ex011
	\item Reformulate the bijection in Theorem~\ref{prop:Yoneda-Ext} as follows.
	From an $n$-extension
	$0\to Y\to E^1\to\cdots\to E^n\to X\to0$, construct
	$t^{-1}b\in\Hom_{\cate{D}(\mathcal{A})}(X,Y[n])$ as follows:
% segment-id: o014.li2.chapter4.k-flat-derived-tensor-b.g006
	\[\begin{tikzcd}
		X \arrow[d, "b"']\arrow[phantom, r, "" {name=A, yshift=2ex}] & {} & & & X \arrow[d, "\identity"] \\
		W & E^1 \arrow[r] & E^2 \arrow[r] & \cdots \arrow[r] & X \\
		Y[n] \arrow[u, "{t: \text{quasi-isomorphism}}"] \arrow[phantom, r, "" {name=B, yshift=-2ex}] & Y \arrow[u] & & &
		\arrow[dash, from=A, to=B]
	\end{tikzcd}\]
	This gives $t^{-1}b\in\Hom_{\cate{D}(\mathcal{A})}(X,Y[n])$. Prove that
	this bijection differs by the factor $(-1)^n$ from the original bijection
	based on $as^{-1}$. In the special case $n=1$, use the $t^{-1}b$ version to
	derive the corresponding version of Proposition~\ref{prop:Ext1-via-ses}
	(the one involving the image of $\identity_Y$).

% segment-id: o014.li2.chapter4.k-flat-derived-tensor-b.hn007
	\begin{hint}
		Prove that $ta+(-1)^{n-1}bs:Z\to W$ is canonically null-homotopic. For
		the second part, note the minus sign in
		Proposition~\ref{prop:ses-vs-triangle}.
	\end{hint}

% segment-id: o014.li2.chapter4.k-flat-derived-tensor-b.ex012
	\item Suppose that all objects $S,T$ in an abelian category $\mathcal{A}$
	satisfy $\Ext^2_{\mathcal{A}}(S,T)=0$. Prove that every
	$X\in\Obj(\cate{D}^{\bdd}(\mathcal{A}))$ is isomorphic to
	$\bigoplus_n\Hm^n(X)[-n]$. Show that $\mathcal{A}:=\Z\dcate{Mod}$
	satisfies this condition.

% segment-id: o014.li2.chapter4.k-flat-derived-tensor-b.hn008
	\begin{hint}
		We may assume that $X$ is a bounded complex. Begin with $n\gg0$. Apply
		Theorem~\ref{prop:Yoneda-Ext} and Proposition~\ref{prop:Yoneda-zero} to
		the $2$-extension
		$0\to\Ker(d^n)\to X^n\xrightarrow{d^n}X^{n+1}\to\Coker(d^n)\to0$,
		then modify $X$ step by step until $d^\bullet=0$, throughout up to
		isomorphism.
	\end{hint}

% segment-id: o014.li2.chapter4.k-flat-derived-tensor-b.ex013
	\item Prove that if $\mathcal{A}$ has enough injective objects,
	$F:\mathcal{A}\to\mathcal{A}'$ is a left-exact additive functor, and
	$\mathrm{R}F$ has finite dimension, then the homomorphism
	$\mathrm{K}_0(\mathcal{A})\simeq\mathrm{K}_0(\cate{D}^{\bdd}(\mathcal{A}))
	\to\mathrm{K}_0(\cate{D}^{\bdd}(\mathcal{A}'))\simeq
	\mathrm{K}_0(\mathcal{A}')$ induced by the triangulated functor
	$\mathrm{R}F$ is determined by
	\[ [A] \mapsto \sum_n (-1)^n \left[\mathrm{R}^n F(A)\right], \quad A \in \Obj(\mathcal{A}). \]

% segment-id: o014.li2.chapter4.k-flat-derived-tensor-b.ex014
	\item Let $R$ be a ring. Prove that a filtered direct limit $\varinjlim$ of
	K-flat complexes in $\cate{C}(R\dcate{Mod})$ is again K-flat.

% segment-id: o014.li2.chapter4.k-flat-derived-tensor-b.ex015
	\item Let $R$ be a ring. If a complex $P$ of $R$-modules is K-flat and
	every $P^n$ is a flat $R$-module, then $P$ is called \emph{strongly
	K-flat}.
% segment-id: o014.li2.chapter4.k-flat-derived-tensor-b.l004
	\begin{enumerate}[(i)]
% segment-id: o014.li2.chapter4.k-flat-derived-tensor-b.i010
		\item Prove that for every $X\in\Obj(\cate{C}(R\dcate{Mod}))$ there is a
		strongly K-flat complex $P$ together with a quasi-isomorphism $P\to X$.
% segment-id: o014.li2.chapter4.k-flat-derived-tensor-b.hn009
		\begin{hint}
			Use the dual of the argument for the existence of K-injective
			resolutions; see \S\ref{sec:K-injectives}.
		\end{hint}
% segment-id: o014.li2.chapter4.k-flat-derived-tensor-b.i011
		\item Write $\mathcal{SF}$ for the full subcategory of
		$\cate{K}(R\dcate{Mod})$ formed by the strongly K-flat complexes.
		Express $\cate{D}(R\dcate{Mod})$ as the localization of $\mathcal{SF}$
		at the quasi-isomorphisms.
	\end{enumerate}

% segment-id: o014.li2.chapter4.k-flat-derived-tensor-b.ex016
	\item (P.\ Berthelot, A.\ Ogus) Let $f$ be a non-zero-divisor in a
	commutative ring $R$. For every $R$-module $M$, write
	$M[f]:=\{m\in M:fm=0\}$. A module $M$ with $M[f]=0$ is called
	\emph{$f$-torsion-free}. First show that an $f$-torsion-free $M$ can be
	embedded as a submodule of
	$M[\frac{1}{f}]=M\dotimes{R}R[\frac{1}{f}]$. Thus $f^mM$ makes sense for
	every $m\in\Z$.

	Let $X$ be a complex of $f$-torsion-free $R$-modules. Using the observation
	above, define a subcomplex $\eta_fX$ of
	$X^\bullet\dotimes{R}R[\frac{1}{f}]$ by
	\[ (\eta_f X)^n := \left\{ x \in f^n X^n : d_X^n(x) \in f^{n+1} X^{n+1} \right\}, \quad n \in \Z . \]
% segment-id: o014.li2.chapter4.k-flat-derived-tensor-b.l005
	\begin{enumerate}[(i)]
% segment-id: o014.li2.chapter4.k-flat-derived-tensor-b.i012
		\item Prove that multiplication degree by degree by $f$ gives an
		isomorphism $\eta_f(X[1])\rightiso(\eta_fX)[1]$.
% segment-id: o014.li2.chapter4.k-flat-derived-tensor-b.i013
		\item Show that if $f,g\in R$ are both non-zero-divisors and every $X^n$
		is $fg$-torsion-free, then $\eta_f\eta_gX=\eta_{fg}X$.
% segment-id: o014.li2.chapter4.k-flat-derived-tensor-b.i014
		\item Give a canonical isomorphism
		$\Hm^n(X)/\Hm^n(X)[f]\rightiso\Hm^n(\eta_fX)$; thus the effect of
		$\eta_f$ is to correct $f$-torsion. Deduce that if $X$ and $Y$ are
		complexes of $f$-torsion-free $R$-modules and $\alpha:X\to Y$ is a
		quasi-isomorphism, then $\alpha$ induces a quasi-isomorphism
		$\eta_f(\alpha):\eta_fX\to\eta_fY$.
% segment-id: o014.li2.chapter4.k-flat-derived-tensor-b.i015
		\item Prove that in this way $\eta_f$ induces an additive functor
		$\mathrm{L}\eta_f:\cate{D}(R\dcate{Mod})\to\cate{D}(R\dcate{Mod})$
		whose restriction to $\mathcal{SF}$ from the preceding exercise equals
		$\eta_f$.
% segment-id: o014.li2.chapter4.k-flat-derived-tensor-b.i016
		\item Show that $\mathrm{L}\eta_f$ is not a triangulated functor. Thus it
		is not a left derived functor in the sense of
		Definition~\ref{def:Verdier-localization-functor}, although it remains a
		right Kan extension.

% segment-id: o014.li2.chapter4.k-flat-derived-tensor-b.hn010
		\begin{hint}
			Take $R=\Z$ and $f=p$ prime. Verify that
			$\mathrm{L}\eta_p(\Z/p\Z)=0$, whereas
			$\mathrm{L}\eta_p(\Z/p^2\Z)\simeq\Z/p\Z$.
		\end{hint}
	\end{enumerate}

% segment-id: o014.li2.chapter4.k-flat-derived-tensor-b.ex017
	\item Continuing the preceding exercise, give a canonical morphism
	$\mathrm{L}\eta_f X \otimesL[R] \mathrm{L}\eta_f Y \to
	\mathrm{L}\eta_f\left( X \otimesL[R] Y \right)$ and show that this
	morphism is symmetric in $X$ and $Y$ in the appropriate sense. In the
	language of \cite[Remark 3.1.8]{Li1}, $\mathrm{L}\eta_f$ is in fact a
	right-lax monoidal functor with respect to $\otimesL[R]$.
\end{Exercises}
